\documentclass{article}

\title{Lab 1 Note}
\author{Ova Zhao}
\date{February 25,2026}



\begin{document}
	\maketitle
	\tableofcontents
	
	\newpage
	
	\section{Introduction}
	
	\LaTeX{} is a high-quality typesetting system; it includes features designed for the production of technical and scientific documentation.\\
	
	In Lab Session 1, we will get started with writing a simple document to	familiarize ourselves with the basic \LaTeX{} working environment. Tutor will guide you in creating your first lab note using \LaTeX{}.\\
	
	Why some paragraphs have tabs but some don't.
	
	\section{Lists}
	
	There are two basic types: ordered/numbered list and unordered/bulleted list.
	
	\subsection{ordered list}
	
	\begin{enumerate}
		\item milk
		\item fruit
		\item meat
	\end{enumerate}
	
	\subsection{unordered lists}
	\begin{itemize}
		\item apple
		\item[-] banana
		\item orange
	\end{itemize}
	
	\subsection{nested lists}
	\begin{enumerate}
		\item milk
		\item fruit
			\begin{itemize}
			\item apple
			\item[-] banana
			\item orange
			\end{itemize}
		\item meat
	\end{enumerate}

	\section{simple math}
	
	One of the main advantages of LATEX is the ease with which mathematical expressions can be written. LATEX provides two writing modes for typesetting mathematics:
	\begin{itemize}
		\item \textbf{inline math} mode used for writing formulas that are part of a paragraph
		\item \textbf{display math} mode used to write expressions that are not part of a text or paragraph and are typeset on separate lines
	\end{itemize}
	
	\subsection{inline math}
	
	$y = \arccos x$ is the number in \([0,\pi]\) for which $\cos y = x$.
	
	\subsection{display math}
	
	$$ -1 \leq \sin \theta \leq 1 $$
	
	\subsection{mixed math}
	We can combine two math modes wisely based on the actual content we want to typeset. For example:\\
	
	The root of the quadratic equation \( ax^2 + bx + c = 0\) is given by \[ x = \frac{-b \pm \sqrt{b^2 - 4ac}}{2a}\] where $\Delta = b^2 - 4ac$ is called the \textit{Discriminant}.
	
	
\end{document}